\section{Zoo Animal Utility}
\label{sec:utility}

This section enumerates the utility functions available to NFT holders in the first drop.

\paragraph{Offers.} NFT owners can switch on the offers feature to allow any user in the Zoo ecosystem to find their NFT and submit an offer on the animal. The owner has an offers page where they can accept, decline, or counter-offer. If an offer is accepted, full ownership of the NFT is transferred to the buyer. The NFT, the collateral locked into it, and its earned value to date are all transferred. The current boost level (as defined by the collateral) remains intact across the transfer, and the original utility persists.

\paragraph{Buy / Sell.} Owners can list NFTs on the Zoo Market for a set price, or place them on auction with or without a floor price. If the listing price is met by a purchaser, ownership transfers on chain along with the NFT plus collateral plus earned value. The current boost level continues unchanged.

\paragraph{Collateralize.} NFT owners can infinitely collateralize their animals for further rewards (\emph{Boosts}) using almost any currency, fiat or crypto. The collateral is locked into the NFT and can only be released by burning the NFT. However, owners may also sell the NFT --- inclusive of its collateral value --- on the market.

\paragraph{Boosts.} Boosts amplify daily rewards. The boost multiplier is determined by collateral added:
\begin{itemize}
  \item \$10--\$99: \textbf{1.1$\times$}
  \item \$100--\$999: \textbf{1.2$\times$}
  \item \$1{,}000--\$9{,}999: \textbf{1.3$\times$}
  \item \$10{,}000+: \textbf{1.4$\times$}
\end{itemize}
Total daily yield is computed as: \emph{daily reward} $+$ (\emph{collateral} $\times$ \emph{boost} $\times$ \emph{days/365}).

\paragraph{Earns.} Each NFT animal has a fixed daily reward. Even if the owner never feeds the animal, it accrues value daily. Adding collateral amplifies the daily reward via the boost mechanism described above. Estimated earning potential over time is visible via the Zoo Reward Calculator.

\paragraph{Create Pool.} A pool is a set price the owner determines for two or more NFTs offered as a collection. Any user can purchase the pool by exchanging the price in any supported currency (fiat or crypto). Payments are converted to \$ZOO and paid out to the seller. Sales are irreversible, but pools can be relisted.

\paragraph{Breeding.} Two adult animals of the same species can be bred to mint an NFT egg of the species' status (Endangered, Rare, or Sublime). There are seven generations allowed; therefore six are breedable.

\paragraph{Hatch Eggs.} Hatching reveals the animal inside an egg, sharing the parents' status. Once hatched, the egg NFT is burned and a new baby animal NFT is born. The baby can be fed crypto or fiat (collateralized) to mint the teen and eventually the adult version.

\paragraph{Burns.} Burning a Zoo NFT returns the rewards earned to date and the collateral locked into it. The Zoo sustainability tax applies (10\% to the DAO, 10\% to the liquidity pool).

\paragraph{Metaverse.} Future versions of Zoo include metaverse land where animals can live. A companion app will let owners place their animals in any physical setting via VR. The NFT file formats (FBX and WebGL) enable upload into any external metaverse landscape.

\section{Market Opportunity}
\label{sec:market}

The market opportunity for Zoo is immense, driven by the convergence of gaming, decentralized finance (DeFi), and the increasing demand for NFTs. The global gaming industry has experienced remarkable growth, reaching a market value of over \$180 billion in 2021, and this upward trend is expected to continue.

Simultaneously, the rise of blockchain technology and DeFi has opened new avenues for innovation, offering exciting possibilities for integrating digital assets and financial systems.

Furthermore, the NFT market has witnessed exponential growth, attracting mainstream attention and transforming various industries. As the NFT space expands, there is a growing appetite for unique and collectible digital assets. Zoo leverages this market demand by creating digital representations of endangered animals as NFTs, providing users with a meaningful and tangible ownership experience while raising awareness about conservation efforts.

The concept of extended reality, including virtual reality (VR) and the metaverse, has gained significant attention from industry leaders and investors. The metaverse represents a potential \$8 trillion to \$13 trillion opportunity by 2030, with a projected user base of up to 5 billion individuals, according to leading financial institutions. Zoo is well-positioned to capitalize on this emerging trend by offering a captivating metaverse environment where users can immerse themselves in a virtual world dedicated to the preservation of endangered species.

Additionally, the integration of play-to-earn and play-and-earn models within the Zoo ecosystem presents an exciting opportunity for users to earn real-world rewards and financial incentives while enjoying the gaming experience. This innovative approach taps into the growing interest in the gamification of finance and the desire for more interactive and rewarding gameplay.

By combining gaming, DeFi, NFTs, and the metaverse, Zoo is poised to capture a significant share of the market. The project's unique value proposition, backed by its commodity-backed native token (\$ZOO) and the emotional intelligence of its digital animal companions, sets it apart in the industry. With its proven track record and the anticipated launch of new game mechanics, staking features, and DAO rewards, Zoo is well-positioned to seize the market opportunity and establish itself as a leading player in the Game-Fi multiverse ecosystem.
